\section{Library Design}\label{sec:design}

This section presents the library bottom-up.  Each subsection names its
home module; code is quoted verbatim.  The design brief throughout: stay on
the standard library's abstractions (setoids, function bundles, algebra
bundles), keep everything \texttt{--safe} and
\texttt{--cubical-compatible}, and arrange definitions so that Agda's
evaluator does the heavy lifting in concrete instances.

\subsection{Words: the free monoid}\label{sec:words}

Module \texttt{Word.Base} defines words over a generator set as an
unquotiented syntax tree --- deliberately \emph{not} lists:

\begin{agdacode}
data Word (X : Set) : Set where
  [_]ʷ : X → Word X
  ε    : Word X
  _•_  : Word X → Word X → Word X
\end{agdacode}

Keeping composition as a binary node means a word remembers its
bracketing, so terms mirror pen-and-paper calculations and rules may be
applied to any parenthesized subterm; associativity is then one of the
\emph{congruence} generators rather than a property of the carrier, and a
reflective solver discharges re-bracketing obligations
(\S\ref{sec:solvers}).  The functorial and monadic structure is standard:
\aid{wmap}, \aid{wconcat}, and the extension of a generator-level
substitution to words, written postfix,

\begin{agdacode}
_ʷ : (X → Word Y) → (Word X → Word Y)
_ʷ = wconcatmap
\end{agdacode}

so \aid{(f ʷ)} is ``the unique monoid homomorphism extending $f$''.  Two
stateful traversals drive everything in \S\ref{sec:engine}: a \emph{coset
action} on letters extends to words by threading the coset left to right,

\begin{agdacode}
_ᵗ : (C → Y → Word X × C) → (C → Word Y → Word X × C)
\end{agdacode}

together with its right-to-left mirror \aid{{\_}ᵗ'}, and conjugation
helpers \aid{{\_}ʰ}, \aid{{\_}ⁿ}, \aid{{\_}ʰ'}, \aid{{\_}ⁿ'} lift an
action of one alphabet on another through words.  A relation set is simply
$\aid{WRel}\;X = \aid{Rel}\;(\aid{Word}\;X)\;0\ell$.

\subsection{Presented monoids as inductive congruences}\label{sec:presentations}

Module \texttt{Presentation.Base}, parameterized by
$\Gamma : \aid{WRel}\;X$, sets the two-level convention used everywhere:
\aid{{\_}==={\_}} \emph{always} means the raw axioms ($\Gamma$ itself), and
\aid{{\_}≈{\_}} \emph{always} means the monoid congruence they generate ---
an inductive family with eight constructors:

\begin{agdacode}
data _≈_ : WRel X where
  refl  : w ≈ w
  sym   : w ≈ v → v ≈ w
  trans : w ≈ v → v ≈ u → w ≈ u
  cong  : w ≈ w' → v ≈ v' → w • v ≈ w' • v'
  assoc      : (w • v) • u ≈ w • (v • u)
  left-unit  : ε • w ≈ w
  right-unit : w • ε ≈ w
  axiom      : w === v → w ≈ v
\end{agdacode}

$\approx$ \emph{is} the word problem of $\langle X \mid \Gamma \rangle$,
and a derivation is a first-class syntactic object --- one can compute
with it, induct on it, and translate it (the congruence-lifting lemmas of
\S\ref{sec:constructions} are inductions on this type).  Compared to
quotient-based designs (e.g.\ mathlib's \aid{PresentedGroup} as a quotient
of a free group by a normal
closure~\cite{mathlib2020,mathlib-doc-presented}, or higher inductive
types in cubical Agda~\cite{vezzosi2019cubical,cubicallibrary}), the
setoid design pays a congruence obligation per function out of the type,
but keeps the normal-form maps of \S\ref{sec:nfproperty} \emph{ordinary
computable functions on words}, with no quotient to eliminate through ---
crucial for the proof-by-evaluation style of the case studies, and
compatible with a standard-library-first policy
(cf.~\cite{barthe2003setoids} for the tradition).
\texttt{Presentation.Properties} packages $\approx$ as an equivalence, a
setoid, and the monoid $\aid{•-ε-monoid}$ --- the presented monoid ---
and \texttt{Presentation.GroupLike} upgrades it to a group when a
\aid{Grouplike} witness supplies a left inverse for every generator:
inverses of words are then derived, cancellation and uniqueness of
inverses are lemmas, and $\aid{•-ε-group}$ is the presented group.

Presentations of concrete algebras are records
(\texttt{Presentation.Definitions}).  The central one:

\begin{agdacode}
record _IsPresentationOf_ (_===_ : WRel X) (G : Group a ℓ) : Set (a ⊔ ℓ) where
  field
    gl : Grouplike _===_
  module GL = Group-Lemmas _===_ gl
  open GroupMorphisms (Group.rawGroup GL.•-ε-group) (Group.rawGroup G)
  field
    ⟦_⟧ : Word X → Group.Carrier G
    iso : IsGroupIsomorphism ⟦_⟧
\end{agdacode}

i.e.\ a grouplike witness plus a group isomorphism from the presented
group onto $G$ whose underlying function is the interpretation.  Weaker
records (\aid{{\_}IsSubPresentationOf{\_}} with a monomorphism;
monoid-level variants) come with upgrade lemmas: a surjective
sub-presentation is a presentation, and a grouplike monoid presentation of
$G$'s underlying monoid is a group presentation --- the latter via the
\texttt{ForStdlib} fact that monoid homomorphisms between groups preserve
inverses (\S\ref{sec:forstdlib}).  These upgrades structure every proof in
the catalogue: one proves soundness, injectivity (by normalization), and
surjectivity separately, and the records assemble the isomorphism.

\subsection{The zoo of normal-form witnesses}\label{sec:nfproperty}

Module \texttt{Normalization.NormalForm.Setoid}, parameterized by $\Gamma$
and a codomain setoid $\mathit{NF}$, defines what it means for a function
out of words to be a normal form.  The definitions are, on the nose, the
standard library's function bundles from the word setoid --- no new
theory, just names for the roles they play:
\[
\begin{array}{lcl}
\aid{NormalFormInjective} & = & \aid{Injection}\ \text{(word setoid)}\ \mathit{NF} \\
\aid{NormalForm}          & = & \aid{RightInverse}\ \text{(word setoid)}\ \mathit{NF} \\
\aid{BijectiveNormalForm} & = & \aid{Bijection}\ \text{(word setoid)}\ \mathit{NF}
\end{array}
\]
An \aid{Injection} bundles the map $\aid{nf}$, its congruence, and
injectivity: $w \approx v \iff \aid{nf}\,w \approx_{\mathit{NF}}
\aid{nf}\,v$.  Two consequences ship with it: the workhorse
\begin{agdacode}
by-equal-nf : nf w ≈ₙ nf v → w ≈ v
\end{agdacode}
which proves word equations by \emph{computing} both normal forms
(\S\ref{sec:solvers}), and decidability transport,
\begin{agdacode}
≈-dec : Decidable _≈ₙ_ → Decidable _≈_
\end{agdacode}
so a normal form with decidable carrier \emph{decides the word problem}.
A \aid{RightInverse} additionally carries a section
$\aid{inv-nf} : |\mathit{NF}| \to \aid{Word}\;X$ with
$\aid{inv-nf} \circ \aid{nf} \approx \mathrm{id}$; from it the library
derives $\aid{normalize} = \aid{inv-nf} \circ \aid{nf}$, its idempotence,
and injectivity of \aid{nf} --- so every \aid{NormalForm} is a
\aid{NormalFormInjective}.

The converse direction is where constructivity bites, a point the library
is careful about.  Classically, any injective $\aid{nf}$ whose image can
be identified has a section by choosing preimages; constructively,
\emph{the section is data}.  An \aid{Injection} gives no algorithm
producing, for a normal form $u$, a word realizing it --- that algorithm
is exactly the extra field of \aid{RightInverse}, and it is what the
completeness lemma below consumes ($\aid{inv-nf}$ realizes each normal
form as a canonical representative).  So the two witnesses are genuinely
different strengths in Agda, injectivity strictly weaker as \emph{data},
even though they classically coincide; the library therefore keeps both,
uses the weak one where only word-problem reasoning is needed
(\aid{by-equal-nf}, decidability), and demands the strong one for
completeness and presentations.  A fourth, deliberately crude witness,
\aid{WeakNormalForm}, drops even the congruence of the map ---
an injective \emph{invariant} suffices to separate words, which is
occasionally all one needs when transporting along unions of relation
sets (\S\ref{sec:constructions}).

Completeness itself is one generic lemma.  Fix a semantics
$\aid{⟦\_⟧} : \aid{Word}\;X \to \mathit{Sem}$ into any setoid.  A
\aid{NormalForm} is \emph{unique for} $\aid{⟦\_⟧}$ when representatives
with equal denotations are equal normal forms:

\begin{agdacode}
record UniqueNormalForm (normalForm : NormalForm) : Set (c ⊔ d) where
  field
    unique : ∀ {u v : |NF|} → ⟦ inv-nf u ⟧ ≈₂ ⟦ inv-nf v ⟧ → u ≈ₙ v

by-normalization : {normalForm : NormalForm} →
  UniqueNormalForm normalForm →
  Congruent _≈_ _≈₂_ ⟦_⟧ → Injective _≈_ _≈₂_ ⟦_⟧
\end{agdacode}

\emph{Soundness plus a unique normal form yields completeness}: given
$\aid{⟦}w\aid{⟧} \approx_2 \aid{⟦}v\aid{⟧}$, soundness transports along
$w \approx \aid{normalize}\,w$, uniqueness identifies the normal forms,
and injectivity of \aid{nf} closes the loop.  The proof is five lines; its
value is that every case study only ever proves \aid{unique} --- a
statement about \emph{normal forms}, i.e.\ about concrete first-order data
--- never about arbitrary words.  A companion lemma
(\aid{by-normalization-wsurj}) similarly reduces surjectivity of the
semantics to surjectivity on normal forms.  Sub-presentations then come
for free (\texttt{Normalization.StarPresentation}): given soundness on
axioms, a \aid{NormalForm}, and a \aid{UniqueNormalForm}, the module
assembles the monoid (or, with \aid{Grouplike}, group) monomorphism ---
the injectivity field being literally \aid{by-normalization}.

\subsection{Inductively defined circuits}\label{sec:circuits}

Module \texttt{Circuit.Base} is parameterized by a gate family
$\aid{Gate} : \mathbb{N} \to \mathsf{Set}$ and defines the wire-indexed
generators:

\begin{agdacode}
data Gen : ℕ → Set where
  gate₁ : Gate 1 → Gen (₁₊ n)
  gate₂ : Gate 2 → Gen (₂₊ n)
  _↥   : Gen n → Gen (₁₊ n)

Circuit : ℕ → Set
Circuit n = Word (Gen n)
\end{agdacode}

A generator is a 1- or 2-ary gate at the bottom of the wires, shifted up
by some number of \aid{{\_}↥}s; $c{\uparrow} = \aid{wmap}\;\aid{{\_}↥}$
shifts a whole circuit.  Two design choices matter more than they look.
First, the indices are built from successors only --- \aid{gate₂} lands in
$\aid{Gen}\,(2{+}n)$, not in $\aid{Gen}\,(k)$ with a side condition
$k \ge 2$, and the $k$-fold shift has type
$\aid{Gen}\;n \to \aid{Gen}\;(k + n)$ with $k$ on the \emph{left} of
$+$.  As the module's comment puts it, this lets the coverage checker
solve its unification problems ``by injectivity of \aid{₁₊} alone,
avoiding stuck Diophantine equations of the form
$\mathit{arity} + k \doteq \mathit{target}$'' --- and downstream, it is
why the coset tables of the case studies can be defined by direct pattern
matching and verified by \aid{refl}.  Second, the structural rules common
to \emph{all} circuit theories are factored out once.  Given any
gate-specific family
$\aid{{\_}SRel,{\_}==={\_}} : (n : \mathbb{N}) \to \aid{CRel}\;n$,

\begin{agdacode}
module Lift-Relation (_SRel,_===_ : (n : ℕ) → CRel n) where
  data _VRel,_===_ : (n : ℕ) → CRel n where
    srel  : n SRel, w === v → n VRel, w === v
    cong↑ : n VRel, w === v → (₁₊ n) VRel, w ↑ === v ↑
    comm₁ : (h : Gate 1) (g : Gen n) → ⋯
    comm₂ : (h : Gate 2) (g : Gen n) → ⋯
\end{agdacode}

adds shift-congruence and far-commutation (a bottom gate commutes with any
generator shifted past its arity; full types in
\S\ref{sec:permutations}).  The accompanying lemma \aid{lemma-cong↑}
lifts the \emph{entire congruence} one wire up --- proved once, by
induction over the eight constructors of $\approx$ --- so equational
developments can move freely along the wire chain
$\mathsf{Cir}\,k \le \mathsf{Cir}\,(k{+}1)$ of
\S\ref{sec:preliminaries}.  Instantiating \aid{Gate} is all it takes to
define a new circuit theory: the permutation instance has one gate and two
axioms; the (in progress) qupit Clifford instance has gates
$S, H, C\!Z, \mathit{Ex}, M$ over a prime parameter.

\subsection{The Reidemeister--Schreier engine}\label{sec:engine}

Module \texttt{Normalization.Reidemeister-Schreier} provides injectivity
theorems for extensions of generator maps, in two strengths.  The
\emph{simplified} form needs no cosets: if $f : X \to \aid{Word}\;Y$ has a
generator-level retraction $g : Y \to \aid{Word}\;X$ ---
$(g\,\aid{ʷ})$ respects $\Delta$'s axioms and
$[x]\aid{ʷ} \approx_1 (g\,\aid{ʷ})(f\,x)$ --- then $(f\,\aid{ʷ})$ is
injective and $(g\,\aid{ʷ})$ surjective.  This is the tool for
\emph{isomorphisms} of presentations, and powers the final step of each
amalgamation case study via the builder \aid{StarIsomorphism}
(\texttt{Presentation.Morphism}), which turns the four generator-level
conditions into an \aid{IsMonoidIsomorphism} of the presented monoids.

The \emph{full} form is coset enumeration.  Its packaging,
\texttt{Normalization.CosetNF}, is the heart of the library
(\aid{SingleLevel} shown; a setoid-coset variant exists):

\begin{agdacode}
module SingleLevel
  (Γ   : WRel X)               -- subgroup presentation  (letters X)
  (Δ   : WRel Y)               -- group presentation     (letters Y)
  (C   : Set)                  -- set of right cosets
  (I   : C)                    -- identity coset
  (f   : X → Word Y)           -- generator embedding  H ↪ G
  (h   : C → Y → Word X × C)   -- coset action (Schreier table)
  ([_] : C → Word Y)           -- Schreier section
\end{agdacode}

Think of $\Gamma$ as presenting a subgroup $H$ of the group $G$ presented
by $\Delta$.  The table $h$ pushes one letter of $G$ past a coset,
emitting a word of $H$ and the successor coset --- so
$\aid{nf} = (h\,\aid{ᵗ})\,I : \aid{Word}\;Y \to \aid{Word}\;X \times C$
runs the whole word from the identity coset.  A submodule
\aid{Transfer} takes the five hypotheses that make the data a genuine
coset presentation --- quoted from the source:

\begin{agdacode}
(h=⁻¹f-gen : ∀ (x : X) → ([ x ]ʷ , I) ~ ((h ᵗ) I (f x)))
(h-wd-ax : ∀ (c : C){u t : Word Y} → u ===₂ t → ((h ᵗ) c u) ~ ((h ᵗ) c t))
(f-wd-ax : ∀ {w v} → w ===₁ v → (f ʷ) w ≈₂ (f ʷ) v)
([I]≈ε : [ I ] ≈₂ ε)
(h=ract :  ∀ c b → let (b' , c') = h c b in
  [ c ] • [ b ]ʷ ≈₂ (f ʷ) b' • [ c' ])
\end{agdacode}

--- (1) the table inverts the embedding at the identity coset; (2) the
table respects $\Delta$'s axioms, coset by coset; (3) the embedding
respects $\Gamma$'s axioms; (4) the identity coset is represented by
$\varepsilon$; (5) sliding one letter past a coset's representative
matches the table.  From these, \aid{Transfer} derives well-definedness of
\aid{nf} on $\approx_2$, the factorization
$[c] \aid{•} b \approx_2 (f\,\aid{ʷ})\,b' \aid{•} [c']$ for whole words,
injectivity --- and the payoff, the \emph{normal-form transports}
\[
  \aid{nfp} : \mathrm{NFI}\;\Gamma\;\mathit{NF}
    \to \mathrm{NFI}\;\Delta\;(\mathit{NF} \times C),
  \qquad
  \aid{nfp'} : \mathrm{NF}\;\Gamma\;\mathit{NF}
    \to \mathrm{NF}\;\Delta\;(\mathit{NF} \times C):
\]
a normal form for the subgroup yields one for the group, with one more
coset digit.  A record \aid{PackedCosetTable} bundles the data and
hypotheses as a single value --- so a whole verified level is an ordinary
first-class object, two of which make an amalgamation
(\S\ref{sec:constructions}) --- and \aid{CosetTower} iterates levels up an
$\mathbb{N}$-indexed family of presentations, with carrier
$\aid{tower-carrier}\;\mathit{NF}_0\;(k{+}1) =
 \aid{tower-carrier}\;\mathit{NF}_0\;k \times C_k$: the mixed-radix normal
form of \S\ref{sec:preliminaries}, one verified digit per level.

Two facts make this engine pleasant in practice.  The hypotheses are
\emph{quantifier-small}: each ranges over cosets and generators (finite
data types) and axioms (a finite inductive family), so in every case
study they are finite case splits.  And the splits close by
\emph{computation}: with $\Gamma$'s normal form in scope, an obligation
like (2) becomes \aid{by-equal-nf Eq.refl} --- both sides normalize to the
same value, and Agda's evaluator checks it.  In the qutrit Clifford+$T$
tower, a nine-coset level imposes seventy-two such obligations; every
one is discharged by \aid{refl}.  This is the precise sense in which the
framework mechanizes the ``gigantic computations'' of
\S\ref{sec:intro}: the computations still happen, but inside the
typechecker.

\subsection{Products of presentations, and their presentation theorems}\label{sec:constructions}

Module \texttt{Presentation.Construct.Base} builds relation sets over a
disjoint union of alphabets from a single primitive, a three-way join
whose \emph{mixed} component is the design parameter:

\begin{agdacode}
data _⋄_⋄_ (Γ₁ : WRel A) (Γ₂ : WRel B) (Γ₃ : WRel (A ⊎ B)) : WRel (A ⊎ B) where
  left  : Γ₁ u v → (Γ₁ ⋄ Γ₂ ⋄ Γ₃) [ u ]ₗ [ v ]ₗ
  right : Γ₂ u v → (Γ₁ ⋄ Γ₂ ⋄ Γ₃) [ u ]ᵣ [ v ]ᵣ
  mid   : Γ₃ u v → (Γ₁ ⋄ Γ₂ ⋄ Γ₃) u v
\end{agdacode}

Choosing the mixed relation recovers the classical constructions ---
free product ($\aid{EmptyRel}$), direct product ($\aid{CommRel}$:
generators of the two sides commute), semidirect product
($\aid{ConjRel}$/$\aid{ConjRelʷ}$: moving $h$ past $n$ conjugates $n$,
by a generator or by a word), amalgamated free product
($\aid{AmalgRel}\,f_1\,f_2$: the two images of a shared alphabet are
identified), and definitional extension ($\aid{SugarRel}$: a new
generator desugars to a word) --- plus $n$-fold iterations
$\Gamma\,\aid{⊕\char94}\,n$.  Congruence-lifting lemmas (\aid{lefts},
\aid{rights}) embed the factors' congruences into the join, and
normal-form transports move witnesses along unions and monomorphisms.

For each construction, \texttt{Presentation.Construct.Properties} proves
two kinds of theorems.  \emph{Normal-form lifting}: from NF witnesses for
$\Gamma$ and $\Delta$, an NF witness for the product relation set, with
product carrier ($\mathit{NF}_1 \times \mathit{NF}_2$ for direct and
semidirect products; for amalgamations the classical alternating carrier
\begin{agdacode}
amalNFC C D = (D ⊎ ⊤) × List (C × D) × (C ⊎ ⊤)
\end{agdacode}
--- a word of alternating nontrivial coset representatives between two
tails~\cite{magnus2004combinatorial}).  \emph{Presentation}: if the factors
present concrete groups, the constructed relation set presents the
constructed group.  Quoting the semidirect and amalgamated statements from
\texttt{MainTheorems.agda}'s aliases (premise telescopes elided):
\[
\begin{array}{l}
\aid{dpres} : (\Gamma \aid{⋄} \Delta \aid{⋄} \aid{CommRel})
   \ \aid{IsPresentationOf}\ (G_1 \times G_2) \\[2pt]
\aid{dpres} : (\Gamma \aid{⋄} \Delta \aid{⋄} \aid{ConjRelʷ}\,\mathit{conj})
   \ \aid{IsPresentationOf}\ (G_1 \rtimes G_2) \\[2pt]
\aid{presentation} : \forall n \to (\Gamma\,\aid{⊕\char94}\,n)
   \ \aid{IsPresentationOf}\ (\otimes\aid{-group}\;n) \\[2pt]
\aid{dpres} : (P_1 \aid{*} P_2 \aid{⋆} f_1 \aid{⋆} f_2)
   \ \aid{IsPresentationOf}\ (G_1 \ast_{G_0} G_2)
\end{array}
\]
where the semantic groups on the right are built by the
\texttt{ForStdlib} constructions below (for the amalgamated product, the
premises are two \aid{PackedCosetTable}s over a common base presentation
--- the record \aid{AmalDataNF} --- plus presentations of the three
groups; the amalgamating homomorphisms $\varphi, \psi$ are then
\emph{derived} from the syntactic embeddings through the presentation
isomorphisms).  A further module proves the recipe for arbitrary group
\emph{extensions}: given a short exact sequence
$1 \to N \to G \to Q \to 1$ realizing conjugation and relator-correction
data, the union of $N$'s relations, word-valued conjugation rules, and
twisted lifts of $Q$'s relators presents $G$
(cf.\ Proposition~2.55 of~\cite{bian2023thesis}); semidirect products are
the special case with trivial corrections.  These construction theorems
are what make the catalogue of \S\ref{sec:examples} cheap: the Pauli
presentation, for instance, is \emph{assembled} --- cyclic factors, one
binary product, one $n$-fold product --- with no fresh coset enumeration
at all.

\subsection{The semantic side: contributions to the standard library}\label{sec:forstdlib}

The right-hand sides of the presentation theorems need the constructed
groups to \emph{exist} as standard-library-style bundles.  The library
ships them in a \texttt{ForStdlib} hierarchy, mirroring stdlib paths for
upstreaming (all \texttt{--safe}, \texttt{--cubical-compatible}):

\begin{itemize}
\item \texttt{Algebra.Construct.SemiDirectProduct}: $N \rtimes H$ for raw
  monoids through groups, driven by a bundled action ---
\begin{agdacode}
record Action (N : RawMonoid a ℓ₁) (H : RawMonoid b ℓ₂) : Set _ where
  field
    act          : H.Carrier → N.Carrier → N.Carrier
    act-cong     : h H.≈ h′ → x N.≈ x′ → act h x N.≈ act h′ x′
    act-ε-homo   : ∀ h → act h N.ε N.≈ N.ε
    act-∙-homo   : ∀ h x y → act h (x N.∙ y) N.≈ act h x N.∙ act h y
    act-identity : ∀ x → act H.ε x N.≈ x
    act-compose  : ∀ h h′ x → act (h H.∙ h′) x N.≈ act h (act h′ x)
\end{agdacode}
  --- ``exactly the laws needed to make the twisted multiplication
  associative and unital'';
\item \texttt{Algebra.Construct.Amalgamation}: the amalgamated free
  product $M \ast_C N$ of \emph{monoids}, with no concrete carrier ---
  elements are words over $|M| \uplus |N|$ and equality is the least
  congruence multiplying out same-factor letters, deleting units, and
  gluing $\varphi\,c$ to $\psi\,c$; when $\varphi, \psi$ are
  homomorphisms this is the pushout in the category of monoids;
\item \texttt{Algebra.Construct.Extension} and
  \texttt{.SplitExtension}: short exact sequences of groups, without and
  with a homomorphic section, semidirect products being the canonical
  split model;
\item \texttt{Algebra.Morphism.Consequences}: a monoid homomorphism
  between groups is a group homomorphism (previously derivable only via a
  deprecated stdlib module);
\item \texttt{Algebra.Morphism.Structures}: \aid{IsMonoidEpimorphism}
  (stdlib has monomorphisms and isomorphisms, but no epimorphisms);
\item \texttt{Data.Fin.Permutation.Properties}: stdlib's
  \aid{Permutation′}~$n$ with composition, identity, and \aid{flip},
  bundled as the \aid{Group} used by the tight semantics of
  \S\ref{sec:permutations}.
\end{itemize}

\subsection{Proof engineering: solvers and evaluation}\label{sec:solvers}

Three small tools carry a large share of the equational text.
\aid{by-assoc} proves $w \approx v$ whenever the two words have the same
letter sequence --- it normalizes both sides to lists and compares them
with \aid{refl} --- so all re-bracketing and unit shuffling in a
calculation is one appeal to the solver; the pattern-guided variant
\aid{by-passoc} re-brackets toward a shape specified by a pattern word
with holes, useful when the \emph{next} rule application needs a
particular association.  \aid{by-equal-nf} (\S\ref{sec:nfproperty})
closes any goal whose two sides have the same computed normal form; since
coset tables and normal-form maps are structurally recursive programs,
this turns table verification into evaluation.  Finally, mundane but
effective: a shared \texttt{Notations} module of numeral patterns
(\aid{₀} … \aid{₁₅}, successor patterns \aid{₁₊} … \aid{₄₊} overloaded
over $\mathbb{N}$ and $\mathsf{Fin}$, and \aid{auto} for \aid{refl})
keeps indices readable, and equational proofs are written in the standard
library's setoid-reasoning style, so the Agda text of, say,
\aid{lemma-XS} in the Clifford+$T$ study reads line for line like the
pen-and-paper derivation it replaces.
